\section{Lemma 3 and Corollary 4}\label{sec:lemma3}

We proved the paper's Lemma 3 and Corollary 4 in Lean, following the paper's argument, and bridged Corollary 4 to OAI's statement of it.

\subsection{Statements}

Lemma3.interpolation, the paper's Lemma 3, is stated at lean-checks/Lemma3.lean:816-825, and that file imports Mathlib only.
The comparator copy lean-checks/comparator-copy/PaperLemma3/Lemma3.lean equals lean-checks/Lemma3.lean through ``end Lemma3'' and has sha256 \hash{1ce196f54fa98689bbea6bd74f2abdb110f5cbcca2eecddcca9316b8af465283}.
It is the file installed for the accepted rerouted Comparator run.
Listing~\ref{lst:interpolation} quotes the statement from the comparator copy, and Listing~\ref{lst:pt} quotes the definition of Pt.
The paper states its Lemma 3 at refs/siegel-paper.tex:252-267.
Its lines 253-266 read:
\begin{paperquote}{3.1}
Let $A:\C^4\to\C^3$ be a surjective linear map, with $\ker A=\C c$, where the
four coordinates of $c$ are linearly independent over $\Q$.
Let $N\ge1$ and $t_1,t_2,t_3$ be integers satisfying
\begin{equation}\label{eq:budget}
 t_j\ge3(N-1)\quad(1\le j\le3),\qquad
 \prod_{j=1}^3\bigl(t_j-3(N-1)+1\bigr)>(4N-3)^4.
\end{equation}
Set $E_N=\{0,\ldots,N-1\}^4$ and
\[
 \cP_t=\{B\in\C[z_1,z_2,z_3]:\deg_{z_j}B\le t_j
                                      \text{ for }1\le j\le3\}.
\]
Then the evaluation map $B\mapsto(B(An))_{n\in E_N}$ from
$\cP_t$ to $\C^{E_N}$ is surjective.
\end{paperquote}
By reading the Lean statement against the paper's text, Lemma3.interpolation matches the paper's Lemma 3, tex 252-267, clause by clause (Read by J. Reilly, 2026-10-27: MATCH).

\begin{listing}[tbp]
\caption{Lemma 3, lean-checks/\allowbreak comparator-copy/\allowbreak PaperLemma3/\allowbreak Lemma3.lean:817-826 (sha256 1ce196f5\ldots).}\label{lst:interpolation}
\begin{Verbatim}
theorem interpolation
    (A : (Fin 4 → ℂ) →ₗ[ℂ] (Fin 3 → ℂ)) (hA : Function.Surjective A)
    (c : Fin 4 → ℂ) (hker : LinearMap.ker A = Submodule.span ℂ {c})
    (hc : LinearIndependent ℚ c)
    (N : ℕ) (hN : 1 ≤ N) (t : Fin 3 → ℤ)
    (ht : ∀ j, 3 * ((N : ℤ) - 1) ≤ t j)
    (hbudget : (4 * (N : ℤ) - 3) ^ 4 < ∏ j, (t j - 3 * ((N : ℤ) - 1) + 1)) :
    Function.Surjective
      (fun B : Pt t => fun n : Fin 4 → Fin N =>
        eval (A (fun i => ((n i : ℕ) : ℂ))) (B : MvPolynomial (Fin 3) ℂ)) := by
\end{Verbatim}
\end{listing}

\begin{listing}[tbp]
\caption{The box Pt, lean-checks/Lemma3.lean:42-43.}\label{lst:pt}
\begin{Verbatim}
def Pt (t : Fin 3 → ℤ) : Set (MvPolynomial (Fin 3) ℂ) :=
  {B | ∀ j, (B.degreeOf j : ℤ) ≤ t j}
\end{Verbatim}
\end{listing}

Lemma3.rectangle, the paper's Corollary 4, is stated in the comparator copy Bridge.lean, installed for the accepted rerouted run (sha256 \hash{af401514808dde5bb913ce8e87cc5280b704a234ea324897d1a1ced1f218f250}).
Listing~\ref{lst:rectangle} quotes it.
The paper states its Corollary 4 at refs/siegel-paper.tex:410-420.
Its lines 411-419 read:
\begin{paperquote}{}
\expandafter\def\csname paperref@lem:interpolation\endcsname{3}\renewcommand{\ref}[1]{\csname paperref@#1\endcsname}
Let $A$ satisfy the hypotheses of Lemma~\ref{lem:interpolation}.
For integers $1\le H\le N$, set
\[
 U=N^{4/3},\qquad T_1=32H^{2/3}U,
 \qquad T_2=T_3=32H^{-1/3}U.
\]
The rows $((An)_1^{\alpha_1}(An)_2^{\alpha_2}(An)_3^{\alpha_3})_{n\in E_N}$,
with $\alpha\in\N^3$ and $\alpha_j\le T_j$, span $\C^{E_N}$.
The constant $32$ is independent of $A$ and $H$.
\end{paperquote}
By reading the Lean statement against the paper's text, Lemma3.rectangle matches the paper's Corollary 4, tex 410-420 (Read by J. Reilly, 2026-10-27: MATCH).
The copies of interpolation, rectangle, bridge\_main, bridge\_descent and via\_lemma3 in the comparator copy and in lean-checks/Lemma3Bridged.lean are identical.

\begin{listing}[tbp]
\caption{Corollary 4, lean-checks/\allowbreak comparator-copy/\allowbreak PaperLemma3/\allowbreak Bridge.lean:96-109 (sha256 af401514\ldots).}\label{lst:rectangle}
\begin{Verbatim}
/-- **Corollary 4** (`cor:rectangle`, tex lines 410–420). -/
theorem rectangle
    (A : (Fin 4 → ℂ) →ₗ[ℂ] (Fin 3 → ℂ)) (hA : Function.Surjective A)
    (c : Fin 4 → ℂ) (hker : LinearMap.ker A = Submodule.span ℂ {c})
    (hc : LinearIndependent ℚ c)
    (H N : ℕ) (hH : 1 ≤ H) (hHN : H ≤ N) :
    Submodule.span ℂ
      ((fun α : Fin 3 → ℕ => fun n : Fin 4 → Fin N =>
          A (fun i => ((n i : ℕ) : ℂ)) 0 ^ α 0 * A (fun i => ((n i : ℕ) : ℂ)) 1 ^ α 1 *
            A (fun i => ((n i : ℕ) : ℂ)) 2 ^ α 2) ''
        {α : Fin 3 → ℕ |
          (α 0 : ℝ) ≤ 32 * (H : ℝ) ^ (2 / 3 : ℝ) * (N : ℝ) ^ (4 / 3 : ℝ) ∧
          (α 1 : ℝ) ≤ 32 * (H : ℝ) ^ (-(1 / 3 : ℝ)) * (N : ℝ) ^ (4 / 3 : ℝ) ∧
          (α 2 : ℝ) ≤ 32 * (H : ℝ) ^ (-(1 / 3 : ℝ)) * (N : ℝ) ^ (4 / 3 : ℝ)}) = ⊤ := by
\end{Verbatim}
\end{listing}

\subsection{Structure of the proof}

Table~\ref{tab:lemma3-structure} lists the sub-lemmas of the Lean proof of Lemma 3 with their line ranges in the comparator copy, grouped under the paper's proof steps (tex 268-408).
The main proof Lemma3.interpolation (816-985) calls each group in order at its ``-- Step'' comments.
Opening the proof, exists\_annihilator and p5\_PtSub give the weights of eq:moments, which is not one of the six named steps.
Six helpers are used across several steps and are not placed in one group: zR 29-30, zC 32-33, emb 35-36, sqd 38-39, Pt 41-43, p5\_eval\_aeval 614-620.

\begin{table}[tbp]
\caption{The Lean proof of Lemma 3 in lean-checks/\allowbreak comparator-copy/\allowbreak PaperLemma3/\allowbreak Lemma3.lean, by the paper's proof steps. Line ranges run from docstring to last line.}\label{tab:lemma3-structure}
\small
\begin{tabularx}{\textwidth}{@{}>{\raggedright\arraybackslash}p{3.3cm}>{\raggedright\arraybackslash}X l@{}}
\toprule
step & sub-lemmas & main proof \\
\midrule
opening, eq:moments (not one of the six steps) & exists\_annihilator 624-661, p5\_PtSub 591-612 & 829-845 \\
dimension count, zeros of R (eq:zeros) & four\_hull\_subset\_cube 663-681, exists\_vanishing 683-738 & 846-885 \\
supporting face and nearest point & exists\_lattice\_nonvanishing 740-770, exists\_nearest 87-148, p1\_var\_ineq 57-85, p1\_coord\_le\_sqd 51-55, p1\_dot\_self\_nonneg 47-49, quarter\_mem\_face 150-208 & 886-909 \\
vertex choice (eq:vertex) & exists\_heavy\_vertex 224-299 & 910-916 \\
off-face vanishing (eq:distance, eq:offface) & closer\_point 301-340, p2\_dot\_expand 213-218, p2\_dot\_self\_nonneg 220-222 & 917-939 \\
Lagrange polynomial Q & exists\_int\_hyperplane 371-432, p3\_g 342-352, p3\_g\_apply 354-354, p3\_clear 356-369, int\_comb\_ne\_zero 434-443, exists\_lagrange 471-586, p4\_linP 445-448, p4\_eval\_linP 450-455, p4\_deg\_linP 457-461, p4\_rL 463-465, p4\_rL\_apply 467-469 & 940-947 \\
final contradiction (B = Q R(z + A(m - u))) & exists\_translate 772-812 & 948-985 \\
\bottomrule
\end{tabularx}
\end{table}

\subsection{The bridge to OAI's Corollary 4}

Lemma3.actual\_biquadratic\_rectangle\_span\_via\_lemma3 is stated in Bridge.lean, and by type\_of\% it has exactly the type of OAI's theorem (Listing~\ref{lst:via}).
The embedding K -> C is chosen as IsAlgClosed.lift, not taken as a hypothesis (Bridge.lean:214).
The kernel vector is OAI's biquadraticCoefficients, defined as \texttt{![1, a⁻¹, -b⁻¹, -(a * b)⁻¹]} at InvariantJetLinearMap.lean:321-322.
ker A = span \{c\} is proved inline in bridge\_main (Bridge.lean:231-254), which is stated at Bridge.lean:196-211.
bridge\_descent (Bridge.lean:147-155) does both the descent from C to K and the restriction to an injective family.
In Mathlib d13f23b, AlgEquiv.mul\_apply gives (e₁ * e₂) x = e₁ (e₂ x), so the third row of A is phi(sigma(tau(1, a, b, ab))).
Listing~\ref{lst:mulapply} quotes its source line, Mathlib/Algebra/Algebra/Equiv.lean:663.
The bridge reuses two OAI setup lemmas: source\_same\_witness\_hyperplane (InvariantJetLinearMap.lean:447-455) and biquadraticCoefficients\_linearIndependent (InvariantJetLinearMap.lean:380-397, with its \_map form at 424-430).

\begin{listing}[tbp]
\caption{The bridge theorem, lean-checks/\allowbreak comparator-copy/\allowbreak PaperLemma3/\allowbreak Bridge.lean:288-290.}\label{lst:via}
\begin{Verbatim}
/-- OAI's Corollary 4, proved via the paper's Lemma 3 instead of the multiplicity estimate. -/
theorem actual_biquadratic_rectangle_span_via_lemma3 :
    type_of% @OAI.SiegelZeros.WeightedTorusJets.actual_biquadratic_rectangle_span := by
\end{Verbatim}
\end{listing}

\begin{listing}[tbp]
\caption{AlgEquiv.mul\_apply, Mathlib d13f23b Mathlib/\allowbreak Algebra/\allowbreak Algebra/\allowbreak Equiv.lean:663, as copied in results/2026-10-26-claims-ledger/mathlib-algequiv-mul\_apply.txt.}\label{lst:mulapply}
\begin{Verbatim}
theorem mul_apply (e₁ e₂ : A₁ ≃ₐ[R] A₁) (x : A₁) : (e₁ * e₂) x = e₁ (e₂ x) :=
\end{Verbatim}
\end{listing}

\subsection{Checks}

All nine declarations compiled from the comparator copy itself (PaperLemma3/Lemma3.lean and Bridge.lean) depend on the axioms [propext, Classical.choice, Quot.sound], at fd4aeeb under the 45-family deny list.
The same axioms were recorded for interpolation, rectangle and via\_lemma3 in Lemma3Bridged.lean at adc7f12 on 2026-10-08.
\footnote{A fresh run on 2026-10-26 at 21:43, of Lemma3Bridged.lean at fd4aeeb under the 45-family deny list, gave the same three axioms lines.
}
A standalone check of the export of via\_lemma3 by nanoda 3a24072 accepted it, with 61810 declarations checked.
The same export with Classical.choice not permitted was rejected (exit 1).
The export with rectangle's proof replaced by interpolation's proof was rejected (exit 101).
These nanoda runs ran as checker via sudo -u checker, with no systemd-run unit and no address-family filter.

An exact execution of the proof of Lemma 3 in Python Fraction arithmetic (src/lemma3\_construct.py) ran 540 instances at N = 2, 3, with 28 checks each, and all pass.
This covers the proof steps at N = 2, 3 only, not the lemma in general.
\footnote{Two working drafts in lean-checks/ are not cited: Lemma3\_Cor4.lean proves interpolation by sorry (line 31), so its rectangle depends on sorryAx, and Lemma3\_Bridge.lean proves rectangle by sorry (line 20).
The canonical Bridge.lean imports neither.
}
